\subsection{Ablation Studies}
\label{sec:exp-ablation}

To isolate the contribution of each design choice we run a series of ablations, each modifying one component of the pipeline while holding the others fixed.

\subsubsection*{(a) Effect of SMOTE on Macro-$F_1$.}
Table~\ref{tab:ablation-smote} reports the change in Macro-$F_1$ when Proportional SMOTE is added on top of the \texttt{class\_weight=balanced} baseline.

\begin{table}[h]
\centering
\caption{$\Delta$Macro-$F_1$ attributable to Proportional SMOTE.}
\label{tab:ablation-smote}
\begin{tabular}{lcccc}
\toprule
\textbf{Model} & \textbf{No SMOTE} & \textbf{With SMOTE} & \textbf{$\Delta$} & \textbf{Verdict} \\
\midrule
LightGBM      & 0.9206 & 0.9241 & $+0.0035$ & \checkmark\ Improved \\
XGBoost       & 0.9214 & 0.9209 & $-0.0005$ & $\approx$ Neutral \\
Random Forest & 0.9202 & 0.9149 & $-0.0053$ & $\times$\ Slightly degraded \\
Decision Tree & 0.9092 & 0.9068 & $-0.0024$ & $\times$\ Slightly degraded \\
\bottomrule
\end{tabular}
\end{table}

\noindent LightGBM is the only model that strictly benefits from SMOTE. Its histogram-based splitting strategy uses new synthetic points to refine bin boundaries in minority-class regions, yielding genuinely better decision surfaces. In contrast, Random Forest already applies \texttt{balanced\_subsample} weighting at every bootstrap draw; adding synthetic neighbours that are close to the originals introduces mild label noise without contributing new information, slightly reducing Macro-$F_1$. XGBoost remains essentially unchanged ($|\Delta| < 0.001$), demonstrating robustness to the oversampling step regardless of whether it helps.

\subsubsection*{(b) Synthetic Fraction vs.\ Reliability.}
Classes with a synthetic fraction above 80\% show degraded reliability in per-class $F_1$. Specifically:
\begin{itemize}
    \item \texttt{Bot} ($\approx 80\%$ synthetic): $F_1 = 0.86$ --- acceptable but not perfect.
    \item \texttt{Web Attack -- XSS} ($\approx 89\%$ synthetic): $F_1 = 0.72$ --- notably degraded.
    \item \texttt{Rare\_Attack} ($\approx 90\%$ synthetic): $F_1 = 0.43$ --- unreliable.
\end{itemize}
This confirms the rule-of-thumb threshold: classes where synthetic samples exceed 80\% of the post-SMOTE total should be treated with caution and their reported metrics annotated accordingly.

\subsubsection*{(c) Consensus of All Four Models (30-sample Qualitative Test).}
As a final qualitative check we randomly drew 30 test flows spanning all 13 classes and obtained independent predictions from all four models. Results:
\begin{itemize}
    \item All four models correct: \textbf{25/30} (83\%)
    \item Majority (3/4) correct: \textbf{0/30} (0\%)
    \item Majority wrong ($\leq$1/4 correct): \textbf{5/30} (17\%)
\end{itemize}
The five failures all belong to the same two hard classes: \texttt{Web Attack -- XSS} (confused unanimously with \texttt{Web Attack -- Brute Force}) and \texttt{Rare\_Attack} (split between \texttt{XSS}, \texttt{BENIGN}, and \texttt{Rare\_Attack}). This inter-model agreement confirms that the confusion is not an artefact of any single learner but is inherent in the flow-level feature space for these attack types.
